\section{Conclusion}
\label{sec:conclusion}

Computer-use agents should be understood as lifecycle-dependent interactive systems rather than benchmark models that merely click on screens. Three takeaways follow. First, CUA capability is lifecycle-dependent: the tri-layer architecture explains how observations become actions, but the four-stage lifecycle explains where that capability is formed, exposed, stressed, and revised. Second, many practical security failures are misattributed as runtime reasoning errors even when their real origin lies in creation-stage data, deployment-stage mediation, or maintenance-stage drift. Third, open deployment settings illustrated by examples such as OpenClaw turn governance into a first-class systems problem, because persistent channels, tool binding, permissions, provenance, and ecosystem evolution jointly determine whether autonomy remains controllable. Progress in CUAs will therefore depend not only on stronger models, but also on safer action abstractions, tighter deployment mediation, better runtime oversight, and continuous maintenance discipline.
